\begin{acknowledgement}
写这致谢的时候，犹豫了半天才动笔。已经好久没有自己写过文章了——AI几乎完全
托管了我的写作能力，连标点符号都差点不知道怎么用。既然是本科的最后一篇
小作文，大抵还是自己写更有价值。为了回顾致谢的写作方式，我读了欧仕刚学长的
本科毕业致谢，整整三页，很难想象AI时代的人类还有能力写如此长的文章。
他是三年前毕业的，里面写了很多新冠疫情时期的感受。时间过得很快，那段日子
对我而言已经如同史前时代，几乎完全忘掉了。

在我本科的这四年，整个世界发生了翻天覆地的变化，其中最令人振奋的无疑是
人工智能的崛起。大一的时候，GPT刚出不久，便已经引发了山呼海啸般的震荡。
许多人相信它将历史性地重塑人类社会，也有些人怀疑这只是又一次的泡沫。
然而海啸愈演愈烈，随着大量资金的涌入与人才的聚集，现在已经没有人能对这头
``房间里的大象''视而不见。如今我们似乎可以下定论：人工智能已经不可阻挡地
涌入这个世界，也终将彻底改变这个世界。从前总有人说，21世纪是生物的世纪
——看起来这预言并没有错，只不过带动生物学发展的并不是生物学家，
而是物理学家与计算机学家。这四年我见证了无数AI圈的大事：OpenAI管理层的
反复内斗，DeepSeek、MiniMax、智谱等国产大模型的横空出世，xAI的诞生与消亡，
Anthropic的暂时辉煌……大量人才涌入这个行业，技术实力与价值观在此激烈碰撞、
对抗，盛况堪比20世纪早期的物理学。我从前一直很悲观，总觉得物理学在21世纪
不会再有大的变化了，能被聪明头脑解决的问题几乎都已被解决完了。
幸好AI的出现让我们又有了更多可以去做的事情。

这篇论文写得并不算符合我的预期。到了北京之后总感觉头脑不是很清醒，
一直达不到尽全力的状态，希望博士入学之后精神状态能更好一些。下面，
正式感谢这段时间里帮助过我的人们。

感谢我的导师王磊老师。这三个月里，他不仅手把手地传授了许多前沿的技术，
更教给我对待学术应有的态度。让我印象最深刻的，是从他身上感受到的一种
``反优绩主义''的氛围——《上交生存手册》里有一句话：``如果一个人把政策评分
作为自己的至高追求，那么他就是这个政策的牺牲品。''在AI急速发展的当下，
许多曾经被视为标准的学术秩序看起来都失去了实际价值。在这种环境下，
真正要做的，是踏踏实实地提高自己的能力、磨砺出属于自己的taste与判断力，
而不是随波逐流地去写一堆没有价值的水文章；要多花时间去思考那些本质上
困难的问题，去解决真正重要的问题。这才是AI时代里一个研究者最值得追求的事，
也是王磊老师在这三个月里反复传递给我的最珍贵的提醒。

感谢这段时间在物理所一同生活、工作的学长们。是你们不计成本、事无巨细的
悉心照顾，帮我解决了所里生活的方方面面——从初到北京的日常起居，到服务器
的使用与配置，再到具体计算问题上的耐心点拨。对一个刚刚进入科研环境的
本科生而言，这种来自学长们的善意，本身就是一种最朴素、也最让人感激的
学术传承。

感谢南京大学物理学院与中国科学院物理研究所的诸位老师。特别感谢我在南京
大学的指导老师吴盛俊教授，在开题阶段和论文撰写过程中给予了我重要的方向
把控；感谢物理所的廖海军老师，在张量网络方法的数值实现与数据分析的关键
节点上给予了我极为耐心的指导与帮助。也感谢这四年来所有曾为我授课、答疑、
讨论过的老师们——你们共同构成了我学术道路上最初的坐标系。

感谢我的朋友管雨泽、聂文辉和潘城坤。我们四个都考到了北京读博，在这几个
月里时常小聚，互相调侃，也互相鼓励，最终都完成了各自的毕业设计。在一个
并不熟悉的城市开启人生新的阶段，能够与你们一同前行，是一件足够幸运的事。

感谢我的父母与家人。这二十多年来，是你们最沉默也最坚定的支持，让我能够
一路沿着自己感兴趣的方向走下去。你们或许并不完全理解我所做的``N皇后''
和``格点气体''究竟是什么，但你们从未质疑过我所走的每一步，永远以最朴素
的方式相信我、托举我。希望从今往后，我也能逐渐成为你们坚实的依靠。

感谢我所生活的国家。这四年里我们也许各自面对过种种小问题，但每当看到
加沙的战火、欧洲的动荡，或是世界其他角落的不安，都会让我无比庆幸——
能在一个没有重大安全问题的中国安心读书、做研究，是一件远比想象中更
珍贵的事。这份太过容易被遗忘的安稳，本身就是这个国家为每一个普通人
默默撑起的一片天。

最后，感谢我所身处的这个时代。它把人工智能、跨学科探索、星辰大海这些
原本只在科幻小说里出现的事物，活生生地摆在了我们这一代人面前。它充满
不确定与焦虑，但同样也意味着前所未有的自由——只要你愿意去做一件真正
想做的事，总能找到属于自己的入口。能在这样的时代开始自己的科研生涯，
是我最大的幸运。愿余下的日子里，我依然有勇气与好奇，去回应这个时代
抛给我的每一个问题。
\end{acknowledgement}
